% GENERATED FILE. Edit canonical lessons, then run npm run generate:books.
% canonical-source-hash: fnv1a64:8fe70adabd312c61
% canonical-lessons: JA-C478-nedaru, JA-C478-amayakasu, JA-C478-ijimeru, JA-C478-karakau, JA-C478-maneru

\chapter{Growing Up Together}
\label{ch:ja-growing-up-together}
\hlchaptermodality{478}

\begin{chapteropening}
\textbf{By the end of this chapter:} \emph{I can talk about children growing up.}

Complete the last lesson of chapter 478: I can talk about children growing up.
\end{chapteropening}

\section[\texorpdfstring{\ja{ねだる} (nedaru)}{ねだる (nedaru)}]{\ja{ねだる} (nedaru) — to pester for; to beg for}

\label{lesson:JA-C478-nedaru}



\begin{quote}
Before the new one: say the Japanese for a bother, then the Japanese for upbringing.
\end{quote}

\subsection*{You'll want to know: \ja{ねだる}}
\textbf{\ja{ねだる}} — \emph{nedaru} — \textquotedblleft{}to pester for; to beg for\textquotedblright{}.

\begin{quote}
\textbf{\ja{すえっこが} \ja{おかしを} \ja{ねだっています。}} — \emph{suekko ga okashi o nedatte imasu} — \textquotedblleft{}The youngest child is pestering for sweets.\textquotedblright{}
\end{quote}

Polite \textbf{\ja{ねだります}}, \emph{nedarimasu}; a godan verb. Past (polite): \textbf{\ja{むすこが} \ja{おかしを} \ja{ねだりました。}}, \emph{musuko ga okashi o nedarimashita}.

\begin{cousinweb}
\textbf{\ja{ねだる}} (\emph{nedaru}) is a native Japanese word (wago), inherited rather than borrowed from Chinese. JMdict marks it as usually written in kana; its kanji spelling borrows characters meaning \textquotedblleft{}strong\textquotedblright{} and \textquotedblleft{}ask\textquotedblright{}, chosen for the meaning and not the sound (JMdict, KANJIDIC2).
\end{cousinweb}

\begin{grammarlens}[title={What you've built}]
One more word for describing what you lived and felt.
\end{grammarlens}

\subsection*{Your turn}
\begin{itemize}
\raggedright
  \item \emph{Say it:} \emph{nedaru}
  \item \emph{Say it:} \emph{nedaru}, once more
  \item \emph{Say it:} \emph{suekko ga okashi o nedatte imasu} — The youngest child is pestering for sweets.
  \item \emph{Recall:} say the Japanese for a bother, then the Japanese for upbringing, then say \emph{nedaru} again
\end{itemize}

\begin{tcolorbox}[breakable,colback=teal!4,colframe=teal!35!black,title={Before you move on}]
What does \ja{ねだる} mean? (\textquotedblleft{}To pester for; to beg for\textquotedblright{}.) Say it once more.
\end{tcolorbox}
\section[\texorpdfstring{\ja{あまやかす} (amayakasu)}{あまやかす (amayakasu)}]{\ja{あまやかす} (amayakasu) — to spoil a child}

\label{lesson:JA-C478-amayakasu}



\begin{quote}
Before the new one: say the Japanese for upbringing, then the Japanese for to pester for.
\end{quote}

\subsection*{You'll want to know: \ja{あまやかす}}
\textbf{\ja{あまやかす}} — \emph{amayakasu} — \textquotedblleft{}to spoil a child\textquotedblright{}.

\begin{quote}
\textbf{\ja{そぼは} \ja{いもうとを} \ja{あまやかしています。}} — \emph{sobo wa imōto o amayakashite imasu} — \textquotedblleft{}My grandmother spoils my younger sister.\textquotedblright{}
\end{quote}

Polite \textbf{\ja{あまやかします}}, \emph{amayakashimasu}; a godan verb. Past (polite): \textbf{\ja{ちちが} \ja{むすめを} \ja{あまやかしました。}}, \emph{chichi ga musume o amayakashimashita}.

\begin{cousinweb}
\textbf{\ja{あまやかす}} (\emph{amayakasu}) is a native Japanese word (wago), inherited rather than borrowed from Chinese; its kanji spelling was attached to it later (JMdict).
\end{cousinweb}

\begin{grammarlens}[title={What you've built}]
One more word for describing what you lived and felt.
\end{grammarlens}

\subsection*{Your turn}
\begin{itemize}
\raggedright
  \item \emph{Say it:} \emph{amayakasu}
  \item \emph{Say it:} \emph{amayakasu}, once more
  \item \emph{Say it:} \emph{sobo wa imōto o amayakashite imasu} — My grandmother spoils my younger sister.
  \item \emph{Recall:} say the Japanese for upbringing, then the Japanese for to pester for, then say \emph{amayakasu} again
\end{itemize}

\begin{tcolorbox}[breakable,colback=teal!4,colframe=teal!35!black,title={Before you move on}]
What does \ja{あまやかす} mean? (\textquotedblleft{}To spoil a child\textquotedblright{}.) Say it once more.
\end{tcolorbox}
\section[\texorpdfstring{\ja{いじめる} (ijimeru)}{いじめる (ijimeru)}]{\ja{いじめる} (ijimeru) — to bully}

\label{lesson:JA-C478-ijimeru}



\begin{quote}
Before the new one: say the Japanese for to pester for, then the Japanese for to spoil a child.
\end{quote}

\subsection*{You'll want to know: \ja{いじめる}}
\textbf{\ja{いじめる}} — \emph{ijimeru} — \textquotedblleft{}to bully\textquotedblright{}.

\begin{quote}
\textbf{\ja{ともだちを} \ja{いじめないで} \ja{ください。}} — \emph{tomodachi o ijimenaide kudasai} — \textquotedblleft{}Please do not bully your friends.\textquotedblright{}
\end{quote}

Polite \textbf{\ja{いじめます}}, \emph{ijimemasu}; an ichidan verb. Past (polite): \textbf{\ja{あにが} \ja{おとうとを} \ja{いじめました。}}, \emph{ani ga otōto o ijimemashita}.

\begin{cousinweb}
\textbf{\ja{いじめる}} (\emph{ijimeru}) is a native Japanese word (wago), inherited rather than borrowed from Chinese; JMdict writes it in kana (JMdict).
\end{cousinweb}

\begin{grammarlens}[title={What you've built}]
One more word for describing what you lived and felt.
\end{grammarlens}

\subsection*{Your turn}
\begin{itemize}
\raggedright
  \item \emph{Say it:} \emph{ijimeru}
  \item \emph{Say it:} \emph{ijimeru}, once more
  \item \emph{Say it:} \emph{tomodachi o ijimenaide kudasai} — Please do not bully your friends.
  \item \emph{Recall:} say the Japanese for to pester for, then the Japanese for to spoil a child, then say \emph{ijimeru} again
\end{itemize}

\begin{tcolorbox}[breakable,colback=teal!4,colframe=teal!35!black,title={Before you move on}]
What does \ja{いじめる} mean? (\textquotedblleft{}To bully\textquotedblright{}.) Say it once more.
\end{tcolorbox}
\section[\texorpdfstring{\ja{からかう} (karakau)}{からかう (karakau)}]{\ja{からかう} (karakau) — to tease}

\label{lesson:JA-C478-karakau}



\begin{quote}
Before the new one: say the Japanese for to spoil a child, then the Japanese for to bully.
\end{quote}

\subsection*{You'll want to know: \ja{からかう}}
\textbf{\ja{からかう}} — \emph{karakau} — \textquotedblleft{}to tease\textquotedblright{}.

\begin{quote}
\textbf{\ja{あには} \ja{いつも} \ja{いもうとを} \ja{からかいます。}} — \emph{ani wa itsumo imōto o karakaimasu} — \textquotedblleft{}My elder brother teases my younger sister all the time.\textquotedblright{}
\end{quote}

Polite \textbf{\ja{からかいます}}, \emph{karakaimasu}; a godan verb. Past (polite): \textbf{\ja{ともだちが} \ja{かれを} \ja{からかいました。}}, \emph{tomodachi ga kare o karakaimashita}.

\begin{cousinweb}
\textbf{\ja{からかう}} (\emph{karakau}) is a native Japanese word (wago), inherited rather than borrowed from Chinese; JMdict writes it in kana (JMdict).
\end{cousinweb}

\begin{grammarlens}[title={What you've built}]
One more word for describing what you lived and felt.
\end{grammarlens}

\subsection*{Your turn}
\begin{itemize}
\raggedright
  \item \emph{Say it:} \emph{karakau}
  \item \emph{Say it:} \emph{karakau}, once more
  \item \emph{Say it:} \emph{ani wa itsumo imōto o karakaimasu} — My elder brother teases my younger sister all the time.
  \item \emph{Recall:} say the Japanese for to spoil a child, then the Japanese for to bully, then say \emph{karakau} again
\end{itemize}

\begin{tcolorbox}[breakable,colback=teal!4,colframe=teal!35!black,title={Before you move on}]
What does \ja{からかう} mean? (\textquotedblleft{}To tease\textquotedblright{}.) Say it once more.
\end{tcolorbox}
\section[\texorpdfstring{\ja{まねる} (maneru)}{まねる (maneru)}]{\ja{まねる} (maneru) — to imitate; to mimic}

\label{lesson:JA-C478-maneru}



\begin{quote}
Before the new one: say the Japanese for to bully, then the Japanese for to tease.
\end{quote}

\subsection*{You'll want to know: \ja{まねる}}
\textbf{\ja{まねる}} — \emph{maneru} — \textquotedblleft{}to imitate; to mimic\textquotedblright{}.

\begin{quote}
\textbf{\ja{こどもは} \ja{ははを} \ja{まねています。}} — \emph{kodomo wa haha o manete imasu} — \textquotedblleft{}The child is imitating her mother.\textquotedblright{}
\end{quote}

Polite \textbf{\ja{まねます}}, \emph{manemasu}; an ichidan verb. Past (polite): \textbf{\ja{おとうとが} \ja{あにを} \ja{まねました。}}, \emph{otōto ga ani o manemashita}.

\begin{cousinweb}
\textbf{\ja{まねる}} (\emph{maneru}) is a native Japanese word (wago), inherited rather than borrowed from Chinese. It shares its first part and its kanji spelling with the noun \emph{mane}, \textquotedblleft{}imitating\textquotedblright{}; that spelling pairs \textquotedblleft{}true\textquotedblright{} and \textquotedblleft{}resemble\textquotedblright{} for the meaning, not the sound (JMdict, KANJIDIC2).
\end{cousinweb}

\begin{grammarlens}[title={What you've built}]
One more word for describing what you lived and felt.
\end{grammarlens}

\subsection*{Your turn}
\begin{itemize}
\raggedright
  \item \emph{Say it:} \emph{maneru}
  \item \emph{Say it:} \emph{maneru}, once more
  \item \emph{Say it:} \emph{kodomo wa haha o manete imasu} — The child is imitating her mother.
  \item \emph{Recall:} say the Japanese for to bully, then the Japanese for to tease, then say \emph{maneru} again
\end{itemize}

\begin{tcolorbox}[breakable,colback=teal!4,colframe=teal!35!black,title={Before you move on}]
What does \ja{まねる} mean? (\textquotedblleft{}To imitate; to mimic\textquotedblright{}.) Say it once more.
\end{tcolorbox}
